\subsection{Constraints on Large Extra Dimensions from Neutrino Oscillation Experiments --- arXiv:hep-ph/0201128}

\textbf{Authors:} H. Davoudiasl, P. Langacker, M. Perelstein

\paragraph{主な主張}
3世代のbulk neutrinoを含む最小LED模型に太陽・大気・原子炉・加速器データを適用し、階層的質量の場合の最も保守的な制限として$R<0.82\,\mu\mathrm{m}$を90\%信頼水準で得る\cite{Davoudiasl:2002LEDConstraints}。

\paragraph{新規性}
KK混合への感度が質量階層とflavourに依存することを利用し、複数の振動実験が余剰次元半径を相補的に制限する枠組みを構築した。

\paragraph{位置付け}
最小LED模型に対する初期の包括的な振動制約であり、後年の実験更新や非最小模型との比較の基準になる。
